\section{引言：LLM 作为科学与数学发现的工具}

本讲由 UT Austin 教授 Swarat Chaudhuri 主讲，重点介绍 LLM 在数学发现和科学发现中的应用，特别强调 LLM 的\textbf{抽象能力}如何加速发现过程。

\begin{knowledgebox}{数学发现与科学发现的流程}
\textbf{数学发现}：建模 $\to$ 猜想 $\to$ 严格推理（证明） $\to$ 反例/新猜想 $\to$ 迭代循环。\\
\textbf{科学发现}：理论构建 $\to$ 假设提出 $\to$ 实验设计 $\to$ 实验验证 $\to$ 新数据 $\to$ 新理论。\\
AI for Math 旨在自动化证明过程；AI for Science 旨在自动化假设生成和实验设计。
\end{knowledgebox}

\subsection{LLM Agent 在发现中的三大优势}

\begin{enumerate}
    \item \textbf{搜索引导}：LLM 内部蕴含大量先验知识，可引导在巨大假设空间中的搜索方向
    \item \textbf{经验学习}：LLM Agent 可与环境交互，从经验中学习（放入 Prompt 或通过 RL 调整权重）
    \item \textbf{抽象能力}：LLM 可以发现和构造新概念与新工具——这是一个全新但极具潜力的方向
\end{enumerate}

